\subsubsection*{Round 31: input views and the byte loop}
Before this round, recovered subkeys are:
\hx{none}.
The clean view returned by \texttt{peel\_sm4} is
\[C_{\rm view}=\hx{681EDF34D206965E86B3E94F536E4246}.\]
The first oracle response for this round and its peeled view are
\[\begin{aligned}F&=\hx{1B50E209D206965E86B3E94F536E4246},\\F_{\rm view}&=\hx{1B50E209D206965E86B3E94F536E4246}.\end{aligned}\]
The filter operates on these views. In the table, ``fresh'' is the key-byte set for this pair and ``remaining'' is the intersection with earlier observations for the same byte.
\begin{center}\small\begin{tabular}{rrrlll}\toprule
Query & Byte & Mask & $\Delta B$ & Fresh & Remaining\\\midrule
1 & 0 & 01 & 4F000000 & 90 91 & 90 91\\
2 & 0 & 02 & 81000000 & 70 72 91 93 & 91\\
3 & 1 & 01 & 00710000 & 24 25 & 24 25\\
4 & 1 & 02 & 00830000 & 24 26 & 24\\
5 & 2 & 01 & 00007900 & A0 A1 & A0 A1\\
6 & 2 & 02 & 00006F00 & A0 A2 & A0\\
7 & 3 & 01 & 00000080 & 12 13 & 12 13\\
8 & 3 & 02 & 000000C9 & 10 12 & 12\\
\bottomrule\end{tabular}\end{center}
Unique bytes are written into the subkey buffer in the following order. Zero-filled positions not yet written are placeholders, not recovered bytes.
\begin{center}\small\begin{tabular}{rrlr}\toprule
Byte written & Value & Buffer after write & Queries used\\\midrule
0 & 91 & 91000000 & 2\\
1 & 24 & 91240000 & 4\\
2 & A0 & 9124A000 & 6\\
3 & 12 & 9124A012 & 8\\
\bottomrule\end{tabular}\end{center}
Only now is \texttt{keys[31]} assigned: $rk_{31}=\hx{9124A012}$. This word may be used to peel the next attacked round.
\subsubsection*{Round 30: input views and the byte loop}
Before this round, recovered subkeys are:
\hx{rk31=9124A012}.
The clean view returned by \texttt{peel\_sm4} is
\[C_{\rm view}=\hx{D206965E86B3E94F536E42467B938F4C}.\]
For the clean ciphertext, the peeling calls perform the following actual calculations in order. The last column is the recovered older state word:
\begin{center}\small\begin{tabular}{rllll}\toprule
Undone round & S-box input & $\tau$ output & $L$ output & Older word\\\midrule
31 & 96FF9D45 & 1A488C73 & 138D5078 & 7B938F4C\\
\bottomrule\end{tabular}\end{center}
The first oracle response for this round and its peeled view are
\[\begin{aligned}F&=\hx{AE21A70A772712DA86B3E94F536E4246},\\F_{\rm view}&=\hx{772712DA86B3E94F536E42467B938F4C}.\end{aligned}\]
The filter operates on these views. In the table, ``fresh'' is the key-byte set for this pair and ``remaining'' is the intersection with earlier observations for the same byte.
\begin{center}\small\begin{tabular}{rrrlll}\toprule
Query & Byte & Mask & $\Delta B$ & Fresh & Remaining\\\midrule
9 & 0 & 01 & 21000000 & 00 01 & 00 01\\
10 & 0 & 02 & 3C000000 & 01 03 & 01\\
11 & 1 & 01 & 00550000 & CE CF & CE CF\\
12 & 1 & 02 & 006D0000 & CD CF & CF\\
13 & 2 & 01 & 00005100 & 72 73 & 72 73\\
14 & 2 & 02 & 0000AB00 & 70 72 & 72\\
15 & 3 & 01 & 000000EB & E4 E5 & E4 E5\\
16 & 3 & 02 & 000000FF & E5 E7 & E5\\
\bottomrule\end{tabular}\end{center}
Unique bytes are written into the subkey buffer in the following order. Zero-filled positions not yet written are placeholders, not recovered bytes.
\begin{center}\small\begin{tabular}{rrlr}\toprule
Byte written & Value & Buffer after write & Queries used\\\midrule
0 & 01 & 01000000 & 10\\
1 & CF & 01CF0000 & 12\\
2 & 72 & 01CF7200 & 14\\
3 & E5 & 01CF72E5 & 16\\
\bottomrule\end{tabular}\end{center}
Only now is \texttt{keys[30]} assigned: $rk_{30}=\hx{01CF72E5}$. This word may be used to peel the next attacked round.
\subsubsection*{Round 29: input views and the byte loop}
Before this round, recovered subkeys are:
\hx{rk31=9124A012, rk30=01CF72E5}.
The clean view returned by \texttt{peel\_sm4} is
\[C_{\rm view}=\hx{86B3E94F536E42467B938F4C893450AD}.\]
For the clean ciphertext, the peeling calls perform the following actual calculations in order. The last column is the recovered older state word:
\begin{center}\small\begin{tabular}{rllll}\toprule
Undone round & S-box input & $\tau$ output & $L$ output & Older word\\\midrule
31 & 96FF9D45 & 1A488C73 & 138D5078 & 7B938F4C\\
30 & AF8156A0 & 6FBFDA1D & 5B32C6F3 & 893450AD\\
\bottomrule\end{tabular}\end{center}
The first oracle response for this round and its peeled view are
\[\begin{aligned}F&=\hx{FD15CC075271D40B4F8E1DBB536E4246},\\F_{\rm view}&=\hx{4F8E1DBB536E42467B938F4C893450AD}.\end{aligned}\]
The filter operates on these views. In the table, ``fresh'' is the key-byte set for this pair and ``remaining'' is the intersection with earlier observations for the same byte.
\begin{center}\small\begin{tabular}{rrrlll}\toprule
Query & Byte & Mask & $\Delta B$ & Fresh & Remaining\\\midrule
17 & 0 & 01 & 3D000000 & 62 63 & 62 63\\
18 & 0 & 02 & 89000000 & 60 62 & 62\\
19 & 1 & 01 & 00E00000 & 28 29 & 28 29\\
20 & 1 & 02 & 001E0000 & 29 2B & 29\\
21 & 2 & 01 & 0000E900 & 34 35 & 34 35\\
22 & 2 & 02 & 00008200 & 34 36 & 34\\
23 & 3 & 01 & 00000057 & 96 97 & 96 97\\
24 & 3 & 02 & 0000001A & 94 96 & 96\\
\bottomrule\end{tabular}\end{center}
Unique bytes are written into the subkey buffer in the following order. Zero-filled positions not yet written are placeholders, not recovered bytes.
\begin{center}\small\begin{tabular}{rrlr}\toprule
Byte written & Value & Buffer after write & Queries used\\\midrule
0 & 62 & 62000000 & 18\\
1 & 29 & 62290000 & 20\\
2 & 34 & 62293400 & 22\\
3 & 96 & 62293496 & 24\\
\bottomrule\end{tabular}\end{center}
Only now is \texttt{keys[29]} assigned: $rk_{29}=\hx{62293496}$. This word may be used to peel the next attacked round.
\subsubsection*{Round 28: input views and the byte loop}
Before this round, recovered subkeys are:
\hx{rk31=9124A012, rk30=01CF72E5, rk29=62293496}.
The clean view returned by \texttt{peel\_sm4} is
\[C_{\rm view}=\hx{536E42467B938F4C893450AD6FE46B75}.\]
For the clean ciphertext, the peeling calls perform the following actual calculations in order. The last column is the recovered older state word:
\begin{center}\small\begin{tabular}{rllll}\toprule
Undone round & S-box input & $\tau$ output & $L$ output & Older word\\\midrule
31 & 96FF9D45 & 1A488C73 & 138D5078 & 7B938F4C\\
30 & AF8156A0 & 6FBFDA1D & 5B32C6F3 & 893450AD\\
29 & C3E0A931 & 928929B3 & E957823A & 6FE46B75\\
\bottomrule\end{tabular}\end{center}
The first oracle response for this round and its peeled view are
\[\begin{aligned}F&=\hx{84EB762070DAF490AA919B05476A5256},\\F_{\rm view}&=\hx{476A52567B938F4C893450AD6FE46B75}.\end{aligned}\]
The filter operates on these views. In the table, ``fresh'' is the key-byte set for this pair and ``remaining'' is the intersection with earlier observations for the same byte.
\begin{center}\small\begin{tabular}{rrrlll}\toprule
Query & Byte & Mask & $\Delta B$ & Fresh & Remaining\\\midrule
25 & 0 & 01 & 04000000 & 42 43 & 42 43\\
26 & 0 & 02 & 55000000 & 40 42 & 42\\
27 & 1 & 01 & 00820000 & 8C 8D & 8C 8D\\
28 & 1 & 02 & 00450000 & 8D 8F & 8D\\
29 & 2 & 01 & 00005800 & 36 37 & 36 37\\
30 & 2 & 02 & 00006000 & 34 36 & 36\\
31 & 3 & 01 & 00000096 & 54 55 & 54 55\\
32 & 3 & 02 & 00000022 & 54 56 & 54\\
\bottomrule\end{tabular}\end{center}
Unique bytes are written into the subkey buffer in the following order. Zero-filled positions not yet written are placeholders, not recovered bytes.
\begin{center}\small\begin{tabular}{rrlr}\toprule
Byte written & Value & Buffer after write & Queries used\\\midrule
0 & 42 & 42000000 & 26\\
1 & 8D & 428D0000 & 28\\
2 & 36 & 428D3600 & 30\\
3 & 54 & 428D3654 & 32\\
\bottomrule\end{tabular}\end{center}
Only now is \texttt{keys[28]} assigned: $rk_{28}=\hx{428D3654}$. The four required subkeys are now known; leave the round loop and reverse the key schedule.
